% ============================================================
%  PEMBAHASAN SUPER DETAIL — LATIHAN SOAL & STUDI KASUS
%  BAB 5: SISTEM PERSAMAAN DAN PERTIDAKSAMAAN LINEAR
%  Catatan: Soal Akhir Bab TIDAK dibahas di sini (sesuai permintaan).
% ============================================================
% ============================================================
%  PREAMBUL BERSAMA — BUKU PEMBAHASAN LATIHAN SOAL
%  Dipakai oleh MAE-2026-2027-Latihan1.tex ... Latihan9.tex
%  Kompilasi: pdfLaTeX (disarankan Overleaf / MiKTeX)
% ============================================================
\documentclass[11pt,a4paper]{report}

% ---------- PAKET ----------
\usepackage[utf8]{inputenc}
\usepackage[T1]{fontenc}
\usepackage[bahasa]{babel}
\usepackage{amsmath,amssymb,amsthm,mathtools}
\usepackage{geometry}
\usepackage{xcolor}
\usepackage{graphicx}
\usepackage{enumitem}
\usepackage{tikz}
\usetikzlibrary{calc,arrows.meta,angles,quotes,positioning,decorations.pathreplacing,patterns}
\usepackage{pgfplots}
\pgfplotsset{compat=1.18}
\usepackage[most]{tcolorbox}
\usepackage{fancyhdr}
\usepackage{titlesec}
\usepackage{array,booktabs}
\usepackage{multicol}
\usepackage{pifont}
\usepackage[colorlinks=true,linkcolor=biru,urlcolor=biru]{hyperref}

\geometry{a4paper,margin=2.5cm,headheight=15pt}

% ---------- WARNA ----------
\definecolor{biru}{HTML}{1F4E79}
\definecolor{birumuda}{HTML}{D6E4F0}
\definecolor{hijau}{HTML}{2E7D32}
\definecolor{hijaumuda}{HTML}{E3F2E1}
\definecolor{oranye}{HTML}{C75B12}
\definecolor{oranyemuda}{HTML}{FBE8DA}
\definecolor{abu}{HTML}{F2F2F2}
\definecolor{ungu}{HTML}{6A1B9A}
\definecolor{ungumuda}{HTML}{EDE2F3}
\definecolor{merah}{HTML}{C62828}
\definecolor{merahmuda}{HTML}{FCE4E4}
\definecolor{tosca}{HTML}{00838F}
\definecolor{toscamuda}{HTML}{E0F2F1}
\definecolor{emas}{HTML}{8D6E00}
\definecolor{emasmuda}{HTML}{FFF6D6}

% ---------- HEADER / FOOTER ----------
\newcommand{\babjudul}{Pembahasan Latihan}
\pagestyle{fancy}
\fancyhf{}
\renewcommand{\headrulewidth}{0.8pt}
\renewcommand{\footrulewidth}{0pt}
\fancyhead[L]{\small\textcolor{ungu}{\leftmark}}
\fancyhead[R]{\small\textcolor{ungu}{\babjudul}}
\fancyfoot[C]{\thepage}

% ---------- GAYA JUDUL BAB ----------
\titleformat{\chapter}[display]
  {\normalfont\huge\bfseries\color{ungu}}
  {\filright\large PEMBAHASAN}{8pt}{\Huge}
\titlespacing*{\chapter}{0pt}{0pt}{20pt}

% ---------- LINGKUNGAN TEOREMA ----------
\newtheoremstyle{gaya}{6pt}{6pt}{}{}{\bfseries\color{biru}}{.}{.5em}{}
\theoremstyle{gaya}
\newtheorem{definisi}{Definisi}[chapter]
\newtheorem{sifat}{Sifat}[chapter]

% ---------- KOTAK: SOAL (ungu) ----------
\newtcolorbox{soalbox}[1]{
  enhanced, breakable, colback=ungumuda, colframe=ungu, boxrule=0.8pt,
  arc=3pt, left=10pt, right=10pt, top=8pt, bottom=8pt,
  title={\textbf{Soal #1}}, fonttitle=\bfseries\color{white}, coltitle=white,
  attach boxed title to top left={xshift=10pt,yshift=-10pt},
  boxed title style={colback=ungu,arc=2pt}}

% ---------- KOTAK: KONSEP KUNCI (hijau) ----------
\newtcolorbox{ingat}[1]{
  enhanced, breakable, colback=hijaumuda, colframe=hijau, boxrule=0.6pt,
  arc=3pt, left=10pt, right=10pt, top=6pt, bottom=6pt,
  title={\textbf{Konsep Kunci: #1}}, fonttitle=\bfseries\color{white}, coltitle=white,
  attach boxed title to top left={xshift=10pt,yshift=-10pt},
  boxed title style={colback=hijau,arc=2pt}}

% ---------- KOTAK: JEBAKAN (merah) ----------
\newtcolorbox{jebakan}{
  enhanced, breakable, colback=merahmuda, colframe=merah, boxrule=0.6pt,
  arc=3pt, left=10pt, right=10pt, top=6pt, bottom=6pt,
  title={\textbf{\ding{55}\ Jebakan \& Kesalahan yang Sering Terjadi}},
  fonttitle=\bfseries\color{white}, coltitle=white,
  attach boxed title to top left={xshift=10pt,yshift=-10pt},
  boxed title style={colback=merah,arc=2pt}}

% ---------- KOTAK: VARIASI SOAL (oranye) ----------
\newtcolorbox{variasi}{
  enhanced, breakable, colback=oranyemuda, colframe=oranye, boxrule=0.6pt,
  arc=3pt, left=10pt, right=10pt, top=6pt, bottom=6pt,
  title={\textbf{\ding{72}\ Bentuk Modifikasi Soal Ini}},
  fonttitle=\bfseries\color{white}, coltitle=white,
  attach boxed title to top left={xshift=10pt,yshift=-10pt},
  boxed title style={colback=oranye,arc=2pt}}

% ---------- KOTAK: STUDI KASUS (biru) ----------
\newtcolorbox{studikasus}[1]{
  enhanced, breakable, colback=abu, colframe=biru, boxrule=0.8pt,
  arc=3pt, left=10pt, right=10pt, top=8pt, bottom=8pt,
  title={\textbf{Studi Kasus: #1}}, fonttitle=\bfseries\color{white}, coltitle=white,
  attach boxed title to top left={xshift=10pt,yshift=-10pt},
  boxed title style={colback=biru,arc=2pt}}

% ---------- KOTAK: TIPS (tosca) ----------
\newtcolorbox{tips}{
  enhanced, breakable, colback=toscamuda, colframe=tosca, boxrule=0.6pt,
  arc=3pt, left=10pt, right=10pt, top=6pt, bottom=6pt,
  borderline west={3pt}{0pt}{tosca}}

% ---------- PERINTAH BANTU ----------
\newcommand{\jawab}{\par\smallskip\noindent\textit{\textcolor{oranye}{Penyelesaian langkah demi langkah:}}\par\smallskip}
\newcommand{\langkah}[1]{\par\smallskip\noindent\textbf{\textcolor{biru}{Langkah #1.}}\ }
\newcommand{\jadi}{\par\smallskip\noindent$\Rightarrow$\ \textbf{Jadi, }}
% Judul tiap nomor soal
\newcommand{\soalno}[2]{%
  \par\bigskip
  \noindent{\large\bfseries\color{ungu}\rule[-2pt]{4pt}{14pt}\ Soal #1}%
  \ {\normalfont\itshape\color{gray}(#2)}\par\nopagebreak\smallskip}

% ---------- HALAMAN JUDUL OTOMATIS ----------
% #1 = judul topik (mis. "Barisan dan Deret"); #2 = nomor bab (angka)
\newcommand{\buathalamanjudul}[2]{%
  \begin{titlepage}
  \centering
  \vspace*{2cm}
  {\color{ungu}\rule{\textwidth}{2pt}}\\[0.6cm]
  {\Huge\bfseries\color{ungu} PEMBAHASAN\\[0.3cm] LATIHAN SOAL}\\[0.4cm]
  {\Large BAB #2 --- #1}\\[0.2cm]
  {\color{ungu}\rule{\textwidth}{2pt}}\\[1.2cm]

  {\large Matematika SMA --- Fase E (Kelas 10)}\\[0.2cm]
  {\large Kurikulum Merdeka}\\[2cm]

  \begin{tcolorbox}[enhanced,width=0.85\textwidth,colback=ungumuda,colframe=ungu,
    arc=4pt,boxrule=1pt]
  \centering\small
  Buku pembahasan ini mengupas \textbf{setiap soal pada kotak ungu (Latihan Soal)}
  di tiap subbab, \emph{selangkah demi selangkah seperti mengajari dari nol}:
  konsep yang dipakai, jebakan yang sering menjerat, cara soal bisa dimodifikasi,
  dan visualisasi. \emph{Soal Akhir Bab tidak dibahas di sini.}
  \end{tcolorbox}

  \vfill
  {\large Tahun Pelajaran 2026/2027}
  \end{titlepage}
  \tableofcontents
  \thispagestyle{empty}
  \clearpage}

% ============================================================
%  PENJAGA KOMPILASI
%  File ini adalah PREAMBUL BERSAMA (di-\input oleh Latihan1..9),
%  BUKAN untuk dikompilasi sendiri. Jika editor terlanjur mengompilasi
%  file ini secara langsung, tampilkan satu halaman info (bukan error).
%  Saat di-\input oleh berkas Latihan, blok ini dilewati.
% ============================================================
\ifnum\pdfstrcmp{\jobname}{MAE-2026-2027-Latihan-preamble}=0
  \begin{document}
  \thispagestyle{empty}
  \begin{center}
  \vspace*{3cm}
  {\Huge\bfseries\color{ungu} Berkas Preambel Bersama}\\[1cm]
  \begin{tcolorbox}[enhanced,width=0.85\textwidth,colback=ungumuda,
    colframe=ungu,arc=4pt,boxrule=1pt]
  \centering
  File ini \textbf{bukan untuk dikompilasi sendiri.} Ia berisi pengaturan
  bersama (paket, warna, kotak, perintah) yang dipanggil lewat
  \texttt{\textbackslash input} oleh berkas \textbf{MAE-2026-2027-Latihan1.tex}
  sampai \textbf{Latihan9.tex}.

  \medskip
  Silakan buka dan kompilasi salah satu berkas \textbf{Latihan1--Latihan9}
  tersebut untuk menghasilkan PDF pembahasan.
  \end{tcolorbox}
  \end{center}
  \end{document}
\fi

\renewcommand{\babjudul}{Pembahasan Latihan — Bab 5}

\begin{document}
\setcounter{chapter}{5}
\buathalamanjudul{Sistem Persamaan dan Pertidaksamaan Linear}{5}

% ============================================================
\chapter*{Pembahasan Bab 5: Sistem Persamaan dan Pertidaksamaan Linear}
\addcontentsline{toc}{chapter}{Pembahasan Bab 5: Sistem Linear}
\markboth{PEMBAHASAN BAB 5}{}
% ============================================================

\begin{tips}
\textbf{Cara membaca buku ini.} Tiap soal: \textbf{(1)} soal ditulis ulang,
\textbf{(2)} konsep kunci, \textbf{(3)} langkah demi langkah, \textbf{(4)}
jebakan, \textbf{(5)} variasi. \emph{Aturan emas Bab 5:} untuk sistem, tujuan
selalu ``menghilangkan satu variabel'' agar masalah mengecil; untuk
pertidaksamaan, tujuannya ``menemukan daerah'', lalu ``memilih titik pojok
terbaik''. \emph{Catatan:} tidak semua sistem berjawaban bilangan bulat ---
pecahan pun jawaban yang sah.
\end{tips}

% ============================================================
\section*{Studi Kasus: Menentukan Harga di Kantin}
\addcontentsline{toc}{section}{Studi Kasus: Harga di Kantin}
% ============================================================

\begin{studikasus}{Menentukan Harga di Kantin}
Tiga teman jajan (harga dalam ribuan): Andi $2$ nasi $+1$ ayam $+1$ es $=32$;
Budi $1$ nasi $+2$ ayam $+1$ es $=34$; Cici $1$ nasi $+1$ ayam $+2$ es $=30$.
\textbf{Pertanyaan.} (a) Bisakah harga ditemukan dari satu baris saja?
(b) Tuliskan secara matematis. (c) Strategi menghilangkan satu variabel?
\end{studikasus}

\subsection*{Jawaban lengkap studi kasus}

\textbf{(a) Mengapa perlu ketiga baris?} Satu persamaan dengan tiga hal tak
diketahui ($x,y,z$) punya tak hingga kemungkinan --- tak bisa memaksa satu harga.
Untuk ``mengunci'' tiga nilai, kita butuh \emph{tiga} informasi bebas.

\textbf{(b) Model matematis (SPLTV).} Misalkan $x=$ harga nasi, $y=$ ayam,
$z=$ es (ribuan):
\[
\begin{cases}
2x+y+z=32\\
x+2y+z=34\\
x+y+2z=30
\end{cases}
\]

\textbf{(c) Strategi: jumlahkan lalu kurangi.}
\langkah{1} Jumlahkan ketiga persamaan. Perhatikan tiap variabel muncul
$2+1+1=4$ kali:
\[
4x+4y+4z=96 \Rightarrow x+y+z=24.
\]
\langkah{2} Kurangkan ``jumlah'' $x+y+z=24$ dari tiap persamaan asli untuk
mengisolasi satu variabel:
\[
(2x+y+z)-(x+y+z)=x=32-24=8,
\]
\[
(x+2y+z)-(x+y+z)=y=34-24=10,\quad
(x+y+2z)-(x+y+z)=z=30-24=6.
\]
\jadi nasi Rp$8.000$, ayam Rp$10.000$, es Rp$6.000$.

\begin{tips}
Trik ``jumlahkan semua lalu kurangi'' bekerja mulus di sini karena koefisiennya
\emph{simetris} (pola $2,1,1$). Untuk sistem umum, pakai eliminasi biasa.
\end{tips}

% ############################################################
\section*{Subbab 5.1 — Sistem Persamaan Linear Tiga Variabel (SPLTV)}
\addcontentsline{toc}{section}{Subbab 5.1 — SPLTV}
% ############################################################

\begin{ingat}{Strategi eliminasi bertingkat}
\textbf{Tujuan:} turunkan 3 variabel $\to$ 2 variabel $\to$ 1 variabel.
\textbf{Langkah:} (1) pilih satu variabel untuk dilenyapkan; (2) pasangkan
persamaan berdua-dua untuk melenyapkannya, hasilnya \emph{dua} persamaan dua
variabel; (3) selesaikan sistem $2\times2$ itu; (4) substitusi balik.
\end{ingat}

% ---------------- SOAL 5.1.1 ----------------
\soalno{5.1.1}{SPLTV eliminasi}
\begin{soalbox}{}
Selesaikan $\begin{cases} 2x+y-z=3\\ x-y+z=2\\ x+y+z=6\end{cases}$.
\end{soalbox}

\jawab
\langkah{1} Beri nama (1),(2),(3). Lenyapkan $z$ dulu.
(2)$+$(3): $(x+x)+(-y+y)+(z+z)=2+6\Rightarrow 2x+2z=8\Rightarrow x+z=4$. \dots(A)
\langkah{2} Lenyapkan $y$ dari (3)$-$(2): $(x-x)+(y+y)+(z-z)=6-2\Rightarrow
2y=4\Rightarrow y=2$.
\langkah{3} Sisipkan $y=2$ ke (1): $2x+2-z=3\Rightarrow 2x-z=1$. \dots(B)
\langkah{4} Selesaikan (A) $x+z=4$ dan (B) $2x-z=1$. Jumlahkan: $3x=5\Rightarrow
x=\frac53$. Maka $z=4-\frac53=\frac73$.
\jadi $(x,y,z)=\left(\dfrac53,\ 2,\ \dfrac73\right)$.
\emph{(Cek pada (1): $2\cdot\frac53+2-\frac73=\frac{10}{3}-\frac73+2=1+2=3$
\checkmark.)}

\begin{jebakan}
\begin{itemize}[leftmargin=*,topsep=2pt]
  \item \textbf{Menolak jawaban pecahan.} SPLTV \emph{tidak wajib} berjawaban
  bulat. Selama lolos pengecekan, pecahan itu benar.
  \item \textbf{Salah tanda saat mengurangi persamaan.} Rapikan: kurangi
  \emph{seluruh} ruas, suku per suku.
\end{itemize}
\end{jebakan}

\begin{variasi}
Ganti satu konstanta agar jawaban bulat; atau minta hanya nilai $x+y+z$
(jumlahkan langsung tanpa mencari tiap variabel).
\end{variasi}

% ---------------- SOAL 5.1.2 ----------------
\soalno{5.1.2}{SPLTV simetris (jumlahkan)}
\begin{soalbox}{}
Selesaikan $\begin{cases} x+y=5\\ y+z=7\\ x+z=6\end{cases}$.
\end{soalbox}

\begin{ingat}{Pola ``berpasangan'' $\Rightarrow$ jumlahkan semua}
Kalau tiap persamaan hanya memuat \emph{dua} variabel dan polanya simetris,
jumlahkan semuanya: tiap variabel muncul dua kali, memberi $2(x+y+z)$.
\end{ingat}

\jawab
\langkah{1} Jumlahkan ketiganya: $2(x+y+z)=5+7+6=18\Rightarrow x+y+z=9$.
\langkah{2} Kurangi tiap persamaan dari total:
\[
z=9-(x+y)=9-5=4,\quad x=9-(y+z)=9-7=2,\quad y=9-(x+z)=9-6=3.
\]
\jadi $(x,y,z)=(2,3,4)$.

\begin{jebakan}
\begin{itemize}[leftmargin=*,topsep=2pt]
  \item \textbf{Lupa membagi 2} setelah menjumlah ($2(x+y+z)=18$, bukan
  $x+y+z=18$).
\end{itemize}
\end{jebakan}

\begin{variasi}
Versi lebih besar (4 variabel berpasangan) memakai ide serupa; atau minta hanya
$x+y+z$.
\end{variasi}

% ---------------- SOAL 5.1.3 ----------------
\soalno{5.1.3}{memodelkan tiga bilangan}
\begin{soalbox}{}
Tiga bilangan berjumlah $30$; yang terbesar $2$ kali terkecil; selisih terbesar
dan tengah adalah $4$. Tentukan ketiganya.
\end{soalbox}

\jawab
\langkah{1} Misalkan terkecil $s$, tengah $m$, terbesar $L$. Terjemahkan:
\[
s+m+L=30,\qquad L=2s,\qquad L-m=4.
\]
\langkah{2} Dari $L=2s$ dan $L-m=4$: $m=L-4=2s-4$.
\langkah{3} Substitusi ke jumlah: $s+(2s-4)+2s=30\Rightarrow 5s-4=30\Rightarrow
5s=34\Rightarrow s=6{,}8$.
\langkah{4} Maka $m=2(6{,}8)-4=9{,}6$ dan $L=2(6{,}8)=13{,}6$.
\jadi ketiga bilangan itu $6{,}8;\ 9{,}6;\ 13{,}6$.
\emph{(Cek: jumlah $=30$; $13{,}6=2\cdot6{,}8$; $13{,}6-9{,}6=4$ \checkmark.)}

\begin{jebakan}
\begin{itemize}[leftmargin=*,topsep=2pt]
  \item \textbf{Menerjemahkan ``2 kali'' terbalik} ($s=2L$). Terbesar $=2\times$
  terkecil berarti $L=2s$.
  \item \textbf{Memaksa bilangan bulat.} Konteks ``bilangan'' di sini boleh
  desimal.
\end{itemize}
\end{jebakan}

\begin{variasi}
Ganti ``selisih $=4$'' agar bulat; atau tambah syarat ketiganya membentuk
barisan aritmetika (kaitan Bab 2).
\end{variasi}

% ---------------- SOAL 5.1.4 ----------------
\soalno{5.1.4}{SPLTV eliminasi campuran}
\begin{soalbox}{}
Selesaikan $\begin{cases} x+2y+z=8\\ 2x-y+z=3\\ x+y-z=0\end{cases}$.
\end{soalbox}

\jawab
\langkah{1} Beri nama (1),(2),(3). Lenyapkan $z$ lewat (2)$+$(3):
$(2x+x)+(-y+y)+(z-z)=3+0\Rightarrow 3x=3\Rightarrow x=1$.
\langkah{2} Lenyapkan $z$ lagi lewat (1)$+$(3): $(x+x)+(2y+y)+(z-z)=8+0\Rightarrow
2x+3y=8$. Masukkan $x=1$: $2+3y=8\Rightarrow y=2$.
\langkah{3} Substitusi ke (3): $1+2-z=0\Rightarrow z=3$.
\jadi $(x,y,z)=(1,2,3)$.
\emph{(Cek (1): $1+4+3=8$ \checkmark.)}

\begin{jebakan}
\begin{itemize}[leftmargin=*,topsep=2pt]
  \item \textbf{Mengeliminasi variabel berbeda tiap pasangan.} Untuk mendapat
  dua persamaan 2-variabel, lenyapkan variabel yang \emph{sama} pada kedua
  pasangan (di sini $z$).
\end{itemize}
\end{jebakan}

\begin{variasi}
Selesaikan dengan substitusi murni (dari (3): $z=x+y$, masukkan ke (1),(2)).
\end{variasi}

% ---------------- SOAL 5.1.5 ----------------
\soalno{5.1.5}{banyak solusi (berimpit)}
\begin{soalbox}{}
Tentukan apakah $\begin{cases} x-y=1\\ 2x-2y=2\end{cases}$ punya satu, tak
hingga, atau tanpa solusi.
\end{soalbox}

\begin{ingat}{Membaca banyaknya solusi}
Bandingkan kedua persamaan: kalau satu adalah \emph{kelipatan} persis yang lain
$\Rightarrow$ \textbf{berimpit} (tak hingga solusi). Kalau ruas kiri kelipatan
tapi ruas kanan tidak $\Rightarrow$ \textbf{sejajar} (tanpa solusi). Selain itu
$\Rightarrow$ satu solusi.
\end{ingat}

\jawab
\langkah{1} Amati: persamaan kedua $2x-2y=2$ adalah persis $2\times$ persamaan
pertama $x-y=1$.
\langkah{2} Karena keduanya persamaan garis yang \emph{sama} (berimpit), setiap
titik pada garis $x-y=1$ memenuhi keduanya.
\jadi sistem punya \textbf{tak hingga} penyelesaian (semua $(x,y)$ dengan
$x-y=1$).

\begin{jebakan}
\begin{itemize}[leftmargin=*,topsep=2pt]
  \item \textbf{Menganggap selalu ada satu jawaban} lalu bingung karena
  variabel ``habis''. Variabel yang lenyap total ($0=0$) menandakan tak hingga
  solusi.
  \item Bandingkan: $2x-2y=5$ (ruas kanan beda) akan memberi $0=3$
  (mustahil) $\Rightarrow$ \emph{tanpa} solusi.
\end{itemize}
\end{jebakan}

\begin{variasi}
Ubah $2x-2y=3$ (sejajar, tanpa solusi) atau $2x-y=2$ (satu solusi). Latih
membedakan ketiga kasus.
\end{variasi}

% ---------------- SOAL 5.1.6 ----------------
\soalno{5.1.6}{SPLDV harga barang}
\begin{soalbox}{}
Harga $2$ buku $+1$ pena $=$ Rp$17$rb; $1$ buku $+3$ pena $=$ Rp$16$rb.
Tentukan harga satuan.
\end{soalbox}

\jawab
\langkah{1} Misalkan $b=$ harga buku, $p=$ harga pena (ribuan):
\[
2b+p=17,\qquad b+3p=16.
\]
\langkah{2} Substitusi: dari persamaan pertama $p=17-2b$. Masukkan ke kedua:
\[
b+3(17-2b)=16\Rightarrow b+51-6b=16\Rightarrow -5b=-35\Rightarrow b=7.
\]
\langkah{3} $p=17-2(7)=3$.
\jadi harga buku Rp$7.000$, pena Rp$3.000$.

\begin{jebakan}
\begin{itemize}[leftmargin=*,topsep=2pt]
  \item \textbf{Salah menuliskan variabel} (tertukar buku/pena). Definisikan
  variabel dengan jelas di awal.
\end{itemize}
\end{jebakan}

\begin{variasi}
Metode eliminasi: kalikan persamaan pertama $3$, kurangi yang kedua untuk
melenyapkan $p$. Jawaban sama.
\end{variasi}

% ---------------- SOAL 5.1.7 ----------------
\soalno{5.1.7}{SPLTV jawaban pecahan}
\begin{soalbox}{}
Selesaikan $\begin{cases} x+y+z=9\\ 2x-y+z=6\\ x+2y-z=3\end{cases}$.
\end{soalbox}

\jawab
\langkah{1} (1)$-$(2): $(x-2x)+(y+y)+(z-z)=9-6\Rightarrow -x+2y=3$. \dots(A)
\langkah{2} (1)$+$(3): $(x+x)+(y+2y)+(z-z)=9+3\Rightarrow 2x+3y=12$. \dots(B)
\langkah{3} Dari (A): $x=2y-3$. Substitusi ke (B): $2(2y-3)+3y=12\Rightarrow
7y-6=12\Rightarrow y=\frac{18}{7}$.
\langkah{4} $x=2\cdot\frac{18}{7}-3=\frac{36-21}{7}=\frac{15}{7}$; lalu
$z=9-x-y=9-\frac{15}{7}-\frac{18}{7}=9-\frac{33}{7}=\frac{30}{7}$.
\jadi $(x,y,z)=\left(\dfrac{15}{7},\ \dfrac{18}{7},\ \dfrac{30}{7}\right)$.
\emph{(Cek (2): $2\cdot\frac{15}{7}-\frac{18}{7}+\frac{30}{7}
=\frac{30-18+30}{7}=\frac{42}{7}=6$ \checkmark.)}

\begin{jebakan}
\begin{itemize}[leftmargin=*,topsep=2pt]
  \item \textbf{Panik melihat pecahan.} Tetap lanjutkan; verifikasi dengan
  substitusi memberi rasa aman.
\end{itemize}
\end{jebakan}

\begin{variasi}
Selesaikan lewat penjumlahan seluruh persamaan untuk mencari kombinasi tertentu,
atau dengan matriks (SNBT/olimpiade).
\end{variasi}

% ---------------- SOAL 5.1.8 ----------------
\soalno{5.1.8}{memodelkan sudut segitiga}
\begin{soalbox}{}
Jumlah tiga sudut segitiga $180^\circ$; sudut $A$ dua kali $B$; $C$ lebih besar
$20^\circ$ dari $A$. Tentukan ketiga sudut.
\end{soalbox}

\jawab
\langkah{1} Terjemahkan: $A+B+C=180$; $A=2B$; $C=A+20$.
\langkah{2} Nyatakan semua dalam $B$: $A=2B$, $C=2B+20$.
\langkah{3} Substitusi ke jumlah:
\[
2B+B+(2B+20)=180\Rightarrow 5B+20=180\Rightarrow 5B=160\Rightarrow B=32^\circ.
\]
\langkah{4} $A=64^\circ$, $C=84^\circ$.
\jadi $A=64^\circ,\ B=32^\circ,\ C=84^\circ$ (jumlah $=180^\circ$ \checkmark).

\begin{jebakan}
\begin{itemize}[leftmargin=*,topsep=2pt]
  \item \textbf{Menerjemahkan ``$C$ lebih besar $20^\circ$ dari $A$'' sebagai
  $C=20-A$.} Yang benar $C=A+20$.
\end{itemize}
\end{jebakan}

\begin{variasi}
Sudut-sudut sebanding $2:3:5$ (bagi $180^\circ$ menurut perbandingan), atau
konteks sudut sepihak/berpelurus.
\end{variasi}

% ---------------- SOAL 5.1.9 ----------------
\soalno{5.1.9}{SPLDV eliminasi}
\begin{soalbox}{}
Selesaikan $\begin{cases} 3x+2y=12\\ x-y=1\end{cases}$.
\end{soalbox}

\jawab
\langkah{1} Dari persamaan kedua: $x=y+1$.
\langkah{2} Substitusi ke pertama: $3(y+1)+2y=12\Rightarrow 5y+3=12\Rightarrow
5y=9\Rightarrow y=\frac95$.
\langkah{3} $x=\frac95+1=\frac{14}{5}$.
\jadi $\left(x,y\right)=\left(\dfrac{14}{5},\ \dfrac95\right)=(2{,}8;\ 1{,}8)$.
\emph{(Cek: $3(2{,}8)+2(1{,}8)=8{,}4+3{,}6=12$ \checkmark.)}

\begin{jebakan}
\begin{itemize}[leftmargin=*,topsep=2pt]
  \item \textbf{Substitusi setengah jalan.} Setelah $x=y+1$, ganti $x$ di
  \emph{seluruh} persamaan pertama, bukan sebagian.
\end{itemize}
\end{jebakan}

\begin{variasi}
Metode eliminasi: kalikan persamaan kedua dengan $2$, jumlahkan agar $y$ lenyap.
\end{variasi}

% ---------------- SOAL 5.1.10 ----------------
\soalno{5.1.10}{model pecahan}
\begin{soalbox}{}
Sebuah pecahan: pembilang $+$ penyebut $=11$; bila pembilang ditambah $3$
nilainya menjadi $1$. Tentukan pecahan itu.
\end{soalbox}

\jawab
\langkah{1} Misalkan pembilang $p$, penyebut $q$. Terjemahkan:
$p+q=11$; dan $\dfrac{p+3}{q}=1\Rightarrow p+3=q$.
\langkah{2} Substitusi $q=p+3$ ke $p+q=11$: $p+(p+3)=11\Rightarrow 2p=8
\Rightarrow p=4$; maka $q=7$.
\jadi pecahannya $\dfrac{4}{7}$.
\emph{(Cek: $4+7=11$; $\frac{4+3}{7}=1$ \checkmark.)}

\begin{jebakan}
\begin{itemize}[leftmargin=*,topsep=2pt]
  \item \textbf{``Nilainya menjadi 1''} berarti pembilang baru $=$ penyebut,
  jadi $p+3=q$ (bukan $p+3=1$).
\end{itemize}
\end{jebakan}

\begin{variasi}
``Bila penyebut dikurangi 2 nilainya $\frac12$'', atau pecahan yang berubah
menjadi $\frac23$ --- semua diterjemahkan ke persamaan.
\end{variasi}

% ############################################################
\section*{Subbab 5.2 — Sistem Pertidaksamaan Linear Dua Variabel (SPtLDV)}
\addcontentsline{toc}{section}{Subbab 5.2 — SPtLDV}
% ############################################################

\begin{ingat}{Menggambar daerah penyelesaian}
(1) Gambar garis batasnya (jadikan tanda ``$=$''). (2) \textbf{Uji satu titik}
(biasanya $(0,0)$): jika memenuhi, arsir sisi yang memuatnya; jika tidak, sisi
seberang. (3) Daerah penyelesaian sistem = \emph{irisan} semua daerah. Kendala
$x\ge0,y\ge0$ mengurung ke kuadran I.
\end{ingat}

% ---------------- SOAL 5.2.1 ----------------
\soalno{5.2.1}{menggambar daerah}
\begin{soalbox}{}
Gambarkan daerah penyelesaian dari $x+2y\le8$, $x\ge0$, $y\ge0$.
\end{soalbox}

\jawab
\langkah{1} Garis batas $x+2y=8$ memotong sumbu di $(8,0)$ (saat $y=0$) dan
$(0,4)$ (saat $x=0$).
\langkah{2} Uji $(0,0)$: $0+0=0\le8$ \emph{benar} $\Rightarrow$ daerah memuat
titik asal (di bawah garis).
\langkah{3} Gabung dengan $x\ge0,y\ge0$ (kuadran I). Hasilnya \emph{segitiga}
bertitik sudut $(0,0),(8,0),(0,4)$.

\begin{center}
\begin{tikzpicture}[scale=0.55]
\fill[hijaumuda](0,0)--(8,0)--(0,4)--cycle;
\draw[->](-0.5,0)--(9.5,0) node[right]{$x$};
\draw[->](0,-0.5)--(0,5.5) node[above]{$y$};
\draw[very thick,hijau](8,0)--(0,4) node[above right]{$x+2y=8$};
\node at (2,1) {\small DP};
\fill (8,0) circle(3pt) node[below]{$(8,0)$};
\fill (0,4) circle(3pt) node[left]{$(0,4)$};
\end{tikzpicture}
\end{center}
\jadi daerahnya segitiga $(0,0),(8,0),(0,4)$.

\begin{jebakan}
\begin{itemize}[leftmargin=*,topsep=2pt]
  \item \textbf{Salah titik potong sumbu.} Untuk $x+2y=8$, titik sumbu-$y$
  adalah $(0,4)$ (bagi $8$ dengan koefisien $2$), bukan $(0,8)$.
  \item \textbf{Lupa uji titik} sehingga mengarsir sisi yang salah.
\end{itemize}
\end{jebakan}

\begin{variasi}
Tanda ``$\ge$'' membalik daerah (menjauhi asal); garis tegak/mendatar
($x\le3$, $y\ge1$).
\end{variasi}

% ---------------- SOAL 5.2.2 ----------------
\soalno{5.2.2}{uji titik}
\begin{soalbox}{}
Tentukan apakah $(1,1)$ memenuhi $2x+3y\le6$.
\end{soalbox}

\jawab
\langkah{1} Substitusi $x=1,y=1$ ke ruas kiri: $2(1)+3(1)=5$.
\langkah{2} Bandingkan: $5\le6$ \emph{benar}.
\jadi \textbf{ya}, $(1,1)$ memenuhi (terletak di daerah penyelesaian).

\begin{jebakan}
\begin{itemize}[leftmargin=*,topsep=2pt]
  \item \textbf{Salah baca tanda.} Kalau tanda ``$\ge$'', maka $5\ge6$
  \emph{salah} --- jawaban berbeda. Selalu perhatikan arah tanda.
\end{itemize}
\end{jebakan}

\begin{variasi}
Uji titik pada beberapa pertidaksamaan sekaligus: titik ada di daerah hanya bila
memenuhi \emph{semua}.
\end{variasi}

% ---------------- SOAL 5.2.3 ----------------
\soalno{5.2.3}{daerah antara dua garis}
\begin{soalbox}{}
Gambarkan $y\ge x$ dan $y\le4$.
\end{soalbox}

\jawab
\langkah{1} $y\ge x$: daerah \emph{di atas} garis $y=x$ (uji $(0,1)$: $1\ge0$
benar).
\langkah{2} $y\le4$: daerah \emph{di bawah} garis mendatar $y=4$.
\langkah{3} Irisannya: pita di atas garis $y=x$ tetapi tidak melebihi $y=4$.

\begin{center}
\begin{tikzpicture}[scale=0.55]
\fill[hijaumuda](-2,-2)--(4,4)--(-2,4)--cycle;
\draw[->](-3,0)--(5,0) node[right]{$x$};
\draw[->](0,-3)--(0,5.5) node[above]{$y$};
\draw[very thick,biru](-2.5,-2.5)--(4.5,4.5) node[right]{$y=x$};
\draw[very thick,oranye](-3,4)--(5,4) node[above right]{$y=4$};
\node at (-0.6,2.6){\small DP};
\end{tikzpicture}
\end{center}
\jadi daerahnya adalah wilayah di atas $y=x$ dan di bawah/pada $y=4$.

\begin{jebakan}
\begin{itemize}[leftmargin=*,topsep=2pt]
  \item \textbf{Menukar atas/bawah garis $y=x$.} $y\ge x$ = titik dengan $y$
  lebih besar = \emph{di atas} garis.
\end{itemize}
\end{jebakan}

\begin{variasi}
Tambah $x\ge0$ untuk membatasi; atau ganti $y\le x$ (daerah di bawah).
\end{variasi}

% ---------------- SOAL 5.2.4 ----------------
\soalno{5.2.4}{menyusun sistem dari daerah}
\begin{soalbox}{}
Tuliskan sistem pertidaksamaan untuk daerah segitiga bertitik $(0,0),(3,0),(0,2)$.
\end{soalbox}

\jawab
\langkah{1} Dua sisi segitiga terletak di sumbu: $x\ge0$ dan $y\ge0$.
\langkah{2} Sisi miring melalui $(3,0)$ dan $(0,2)$. Persamaannya (bentuk
intersep) $\dfrac x3+\dfrac y2=1$, atau dikali $6$: $2x+3y=6$.
\langkah{3} Segitiga berada di sisi yang memuat $(0,0)$: uji $2(0)+3(0)=0\le6$
benar, jadi $2x+3y\le6$.
\jadi sistemnya: $\boxed{x\ge0,\ y\ge0,\ 2x+3y\le6}$.

\begin{jebakan}
\begin{itemize}[leftmargin=*,topsep=2pt]
  \item \textbf{Salah menyusun garis sisi miring.} Gunakan bentuk intersep
  $\frac{x}{a}+\frac{y}{b}=1$ dengan $a,b$ perpotongan sumbu.
\end{itemize}
\end{jebakan}

\begin{variasi}
Segitiga dengan salah satu titik bukan di sumbu (butuh dua garis miring); atau
daerah segi empat (empat pertidaksamaan).
\end{variasi}

% ---------------- SOAL 5.2.5 ----------------
\soalno{5.2.5}{maksimum fungsi objektif}
\begin{soalbox}{}
Maksimumkan $f=x+y$ pada $x+2y\le6$, $2x+y\le6$, $x,y\ge0$.
\end{soalbox}

\begin{ingat}{Teorema titik pojok}
Nilai maksimum/minimum fungsi linear $f=ax+by$ pada daerah tertutup selalu di
salah satu \emph{titik pojok}. Jadi: cari semua titik pojok, hitung $f$ di
masing-masing, pilih yang terbesar/terkecil.
\end{ingat}

\jawab
\langkah{1} Titik pojok: $(0,0)$; $(3,0)$ (dari $2x+y=6$ saat $y=0$); $(0,3)$
(dari $x+2y=6$ saat $x=0$); dan perpotongan $x+2y=6$ dengan $2x+y=6$.
\langkah{2} Perpotongan: kurangkan $(x+2y)-(2x+y)=0\Rightarrow -x+y=0\Rightarrow
x=y$. Lalu $3x=6\Rightarrow x=y=2$, titik $(2,2)$.
\langkah{3} Hitung $f=x+y$ di tiap pojok:
\[
f(0,0)=0,\quad f(3,0)=3,\quad f(0,3)=3,\quad f(2,2)=4.
\]
\jadi maksimum $f=4$ di titik $(2,2)$.

\begin{jebakan}
\begin{itemize}[leftmargin=*,topsep=2pt]
  \item \textbf{Lupa titik pojok perpotongan dua kendala miring.} Sering titik
  optimum justru di sana, bukan di sumbu.
\end{itemize}
\end{jebakan}

\begin{variasi}
Ganti fungsi objektif ($f=3x+y$) menggeser letak optimum; tambah kendala.
\end{variasi}

% ---------------- SOAL 5.2.6 ----------------
\soalno{5.2.6}{minimum pada daerah tak terbatas}
\begin{soalbox}{}
Minimumkan $f=2x+3y$ pada $x+y\ge4$, $x,y\ge0$.
\end{soalbox}

\jawab
\langkah{1} Daerah ($x+y\ge4$ di kuadran I) \emph{tak terbatas} ke atas, tetapi
untuk \emph{minimum} kita hanya perlu titik pojok terbawah.
\langkah{2} Titik pojok pada garis $x+y=4$: $(4,0)$ dan $(0,4)$.
\langkah{3} Hitung: $f(4,0)=2(4)+3(0)=8$; $f(0,4)=2(0)+3(4)=12$.
\jadi minimum $f=8$ di titik $(4,0)$.

\begin{jebakan}
\begin{itemize}[leftmargin=*,topsep=2pt]
  \item \textbf{Mencari maksimum di daerah tak terbatas.} Untuk arah ``$\ge$''
  yang membuka ke atas, \emph{maksimum tidak ada}; minimum tetap di titik pojok.
\end{itemize}
\end{jebakan}

\begin{variasi}
Fungsi $f=3x+2y$ memindah minimum ke $(0,4)$? Cek: $f(4,0)=12$, $f(0,4)=8$
$\Rightarrow$ minimum di $(0,4)$. Arah objektif menentukan pojok.
\end{variasi}

% ---------------- SOAL 5.2.7 ----------------
\soalno{5.2.7}{daerah segi empat}
\begin{soalbox}{}
Gambarkan daerah $x+y\le5$, $x\le3$, $y\ge0$, $x\ge0$.
\end{soalbox}

\jawab
\langkah{1} Batas: garis $x+y=5$ (di bawahnya), garis tegak $x=3$ (di kirinya),
sumbu-sumbu ($x,y\ge0$).
\langkah{2} Cari titik pojok: $(0,0)$; $(3,0)$; perpotongan $x=3$ dengan
$x+y=5$ yaitu $(3,2)$; dan $(0,5)$.
\jadi daerahnya segi empat bertitik $(0,0),(3,0),(3,2),(0,5)$.

\begin{center}
\begin{tikzpicture}[scale=0.55]
\fill[hijaumuda](0,0)--(3,0)--(3,2)--(0,5)--cycle;
\draw[->](-0.5,0)--(6,0) node[right]{$x$};
\draw[->](0,-0.5)--(0,6.5) node[above]{$y$};
\draw[very thick,hijau](5,0)--(0,5);
\draw[very thick,oranye](3,-0.3)--(3,3) node[above]{$x=3$};
\node at (1.2,1.4){\small DP};
\fill (3,2) circle(3pt) node[right]{$(3,2)$};
\fill (0,5) circle(3pt) node[left]{$(0,5)$};
\end{tikzpicture}
\end{center}

\begin{jebakan}
\begin{itemize}[leftmargin=*,topsep=2pt]
  \item \textbf{Lupa memotong dengan $x=3$}, sehingga menganggap segitiga
  $(0,0),(5,0),(0,5)$. Garis tegak ``memotong'' pojok kanan.
\end{itemize}
\end{jebakan}

\begin{variasi}
Tambah $y\le4$ untuk memotong pojok atas; daerah jadi segi lima.
\end{variasi}

% ---------------- SOAL 5.2.8 ----------------
\soalno{5.2.8}{memodelkan kendala}
\begin{soalbox}{}
Sebuah pabrik membuat $2$ produk. Tuliskan kendala jika bahan A tersedia
maksimum $10$ unit dan bahan B maksimum $8$ unit.
\end{soalbox}

\jawab
\langkah{1} Misalkan $x=$ banyak produk 1, $y=$ banyak produk 2. Keduanya tak
mungkin negatif: $x\ge0,\ y\ge0$.
\langkah{2} Nyatakan pemakaian bahan. Misalkan tiap produk 1 memakai $a_1$ unit A
dan $b_1$ unit B; produk 2 memakai $a_2$ unit A dan $b_2$ unit B. Maka:
\[
a_1x+a_2y\le10\ (\text{bahan A}),\qquad b_1x+b_2y\le8\ (\text{bahan B}).
\]
\jadi model umum: $x\ge0,\ y\ge0,\ a_1x+a_2y\le10,\ b_1x+b_2y\le8$.
(Angka $a_i,b_i$ diisi sesuai data pemakaian per produk.)

\begin{jebakan}
\begin{itemize}[leftmargin=*,topsep=2pt]
  \item \textbf{Lupa $x\ge0,y\ge0$.} Jumlah produk tak boleh negatif --- kendala
  ini wajib pada model nyata.
  \item \textbf{Memakai ``$=$'' bukan ``$\le$''.} ``Maksimum'' berarti boleh
  kurang, jadi $\le$.
\end{itemize}
\end{jebakan}

\begin{variasi}
Jika ada batas \emph{minimum} permintaan ($x\ge2$), tambahkan sebagai kendala
``$\ge$''.
\end{variasi}

% ---------------- SOAL 5.2.9 ----------------
\soalno{5.2.9}{cek titik dalam daerah}
\begin{soalbox}{}
Apakah titik $(2,2)$ ada dalam daerah $x+y\le3$?
\end{soalbox}

\jawab
\langkah{1} Substitusi: $2+2=4$.
\langkah{2} Bandingkan: $4\le3$ \emph{salah}.
\jadi \textbf{tidak}, $(2,2)$ berada di \emph{luar} daerah.

\begin{jebakan}
\begin{itemize}[leftmargin=*,topsep=2pt]
  \item \textbf{Menebak dari gambar kasar.} Substitusi angka lebih pasti daripada
  menerka posisi.
\end{itemize}
\end{jebakan}

\begin{variasi}
$(1,1)$ ($=2\le3$ benar, di dalam); titik tepat di garis $(1,2)$ ($=3$, di
\emph{batas} --- termasuk karena tanda ``$\le$'').
\end{variasi}

% ---------------- SOAL 5.2.10 ----------------
\soalno{5.2.10}{menentukan titik pojok}
\begin{soalbox}{}
Tentukan titik pojok dari $x+y\le4$, $x\ge1$, $y\ge0$.
\end{soalbox}

\jawab
\langkah{1} Batas: $x+y=4$, garis tegak $x=1$, sumbu $y=0$.
\langkah{2} Cari perpotongan pasangan garis:
\begin{itemize}[leftmargin=*,topsep=1pt]
  \item $x=1$ \& $y=0$ $\Rightarrow (1,0)$;
  \item $x+y=4$ \& $y=0$ $\Rightarrow (4,0)$;
  \item $x+y=4$ \& $x=1$ $\Rightarrow (1,3)$.
\end{itemize}
\jadi titik pojoknya $(1,0),\ (4,0),\ (1,3)$.

\begin{jebakan}
\begin{itemize}[leftmargin=*,topsep=2pt]
  \item \textbf{Memasukkan $(0,0)$.} Kendala $x\ge1$ menyingkirkan titik asal;
  pojok kiri-bawah bergeser ke $(1,0)$.
\end{itemize}
\end{jebakan}

\begin{variasi}
Tambah $y\le2$: pojok $(1,3)$ terpotong menjadi dua pojok baru.
\end{variasi}

% ############################################################
\section*{Subbab 5.3 — Fungsi, Persamaan, dan Pertidaksamaan Linear (Pengayaan)}
\addcontentsline{toc}{section}{Subbab 5.3 — Fungsi \& Garis Linear}
% ############################################################

\begin{ingat}{Garis lurus $y=mx+c$}
$m=\dfrac{\Delta y}{\Delta x}$ (gradien: ``naik dibagi geser''); $c$ = nilai saat
$x=0$. Sejajar $\iff m_1=m_2$; tegak lurus $\iff m_1 m_2=-1$. Pertidaksamaan
satu variabel: kali/bagi bilangan \emph{negatif} $\Rightarrow$ balik tanda.
\end{ingat}

% ---------------- SOAL 5.3.1 ----------------
\soalno{5.3.1}{gradien dari dua titik}
\begin{soalbox}{}
Tentukan gradien garis melalui $(2,3)$ dan $(6,11)$.
\end{soalbox}

\jawab
\langkah{1} $m=\dfrac{y_2-y_1}{x_2-x_1}=\dfrac{11-3}{6-2}=\dfrac{8}{4}=2$.
\jadi gradiennya $2$ (tiap $x$ naik $1$, $y$ naik $2$).

\begin{jebakan}
\begin{itemize}[leftmargin=*,topsep=2pt]
  \item \textbf{Membalik $\Delta x$ dan $\Delta y$.} Gradien = perubahan
  \emph{$y$} dibagi perubahan \emph{$x$}, bukan sebaliknya.
  \item \textbf{Urutan titik tidak konsisten.} Pastikan $y_2,x_2$ dari titik yang
  sama.
\end{itemize}
\end{jebakan}

\begin{variasi}
Gradien negatif (garis turun); gradien $0$ (mendatar); dua titik dengan $x$ sama
(garis tegak, gradien tak terdefinisi).
\end{variasi}

% ---------------- SOAL 5.3.2 ----------------
\soalno{5.3.2}{persamaan garis titik-gradien}
\begin{soalbox}{}
Susun persamaan garis bergradien $-2$ melalui $(1,4)$.
\end{soalbox}

\jawab
\langkah{1} Bentuk titik-gradien: $y-y_1=m(x-x_1)$ dengan $m=-2$, $(x_1,y_1)=(1,4)$:
\[
y-4=-2(x-1).
\]
\langkah{2} Rapikan: $y-4=-2x+2\Rightarrow y=-2x+6$.
\jadi $y=-2x+6$.

\begin{jebakan}
\begin{itemize}[leftmargin=*,topsep=2pt]
  \item \textbf{Salah distribusi tanda.} $-2(x-1)=-2x+2$ (bukan $-2x-2$).
\end{itemize}
\end{jebakan}

\begin{variasi}
Diberi dua titik (cari $m$ dulu); atau garis melalui titik dengan syarat
sejajar/tegak lurus garis lain.
\end{variasi}

% ---------------- SOAL 5.3.3 ----------------
\soalno{5.3.3}{garis melalui titik asal}
\begin{soalbox}{}
Tentukan persamaan garis melalui $(0,0)$ dan $(3,6)$.
\end{soalbox}

\jawab
\langkah{1} Gradien: $m=\dfrac{6-0}{3-0}=2$.
\langkah{2} Karena melalui $(0,0)$, intersep $c=0$: $y=2x$.
\jadi $y=2x$.

\begin{jebakan}
\begin{itemize}[leftmargin=*,topsep=2pt]
  \item \textbf{Menambah konstanta.} Garis lewat titik asal $\Rightarrow c=0$;
  tak ada ``$+c$''.
\end{itemize}
\end{jebakan}

\begin{variasi}
Garis lewat asal dengan gradien pecahan; atau garis $y=kx$ yang melalui titik
tertentu (cari $k$).
\end{variasi}

% ---------------- SOAL 5.3.4 ----------------
\soalno{5.3.4}{uji sejajar}
\begin{soalbox}{}
Apakah $y=3x+1$ dan $y=3x-4$ sejajar? Jelaskan.
\end{soalbox}

\jawab
\langkah{1} Bandingkan gradien: keduanya $m=3$ (sama), tetapi intersep berbeda
($1$ vs $-4$).
\jadi \textbf{ya, sejajar} --- gradien sama membuat kemiringannya identik,
sedangkan intersep berbeda memastikan keduanya \emph{tidak berimpit} (tak pernah
berpotongan).

\begin{jebakan}
\begin{itemize}[leftmargin=*,topsep=2pt]
  \item \textbf{Menganggap sama = berimpit.} Sejajar butuh gradien sama
  \emph{dan} intersep berbeda; kalau intersep juga sama, keduanya garis yang
  sama.
\end{itemize}
\end{jebakan}

\begin{variasi}
$y=3x+1$ dan $6x-2y=-2$ (sederhanakan: $y=3x+1$, ternyata \emph{berimpit}).
\end{variasi}

% ---------------- SOAL 5.3.5 ----------------
\soalno{5.3.5}{gradien tegak lurus}
\begin{soalbox}{}
Tentukan gradien garis yang tegak lurus $y=\tfrac14 x+2$.
\end{soalbox}

\begin{ingat}{Syarat tegak lurus}
$m_1\cdot m_2=-1$, jadi $m_2=-\dfrac{1}{m_1}$ (kebalikan negatif). Bandingkan uji
tegak lurus vektor $\vec a\cdot\vec b=0$ di Bab 3.
\end{ingat}

\jawab
\langkah{1} Gradien acuan $m_1=\dfrac14$.
\langkah{2} $m_2=-\dfrac{1}{m_1}=-\dfrac{1}{1/4}=-4$.
\jadi gradien tegak lurusnya $-4$.

\begin{jebakan}
\begin{itemize}[leftmargin=*,topsep=2pt]
  \item \textbf{Hanya membalik tanda} ($-\frac14$) atau hanya membalik pecahan
  ($4$). Harus \emph{keduanya}: kebalikan \textbf{dan} negatif.
\end{itemize}
\end{jebakan}

\begin{variasi}
Tegak lurus terhadap $y=-2x+5$ $\Rightarrow m_2=\frac12$; garis mendatar tegak
lurus garis tegak.
\end{variasi}

% ---------------- SOAL 5.3.6 ----------------
\soalno{5.3.6}{domain \& range terbatas}
\begin{soalbox}{}
Tentukan domain dan range $f(x)=2x-1$ untuk $-1\le x\le3$.
\end{soalbox}

\jawab
\langkah{1} Domain langsung dari batas: $\{x\mid -1\le x\le3\}$.
\langkah{2} Karena fungsi \emph{naik} ($m=2>0$), nilai terkecil di $x=-1$ dan
terbesar di $x=3$:
\[
f(-1)=2(-1)-1=-3,\qquad f(3)=2(3)-1=5.
\]
\jadi domain $[-1,3]$, range $[-3,5]$.

\begin{jebakan}
\begin{itemize}[leftmargin=*,topsep=2pt]
  \item \textbf{Menukar ujung range} kalau fungsi \emph{turun} ($m<0$): nilai
  terbesar justru di $x$ terkecil. Cek tanda $m$.
\end{itemize}
\end{jebakan}

\begin{variasi}
$f(x)=-3x+2$ pada $[0,4]$ (turun, range terbalik). Domain kontekstual (jumlah
barang bilangan cacah).
\end{variasi}

% ---------------- SOAL 5.3.7 ----------------
\soalno{5.3.7}{persamaan linear satu variabel}
\begin{soalbox}{}
Selesaikan $4x-7=2x+5$.
\end{soalbox}

\jawab
\langkah{1} Kumpulkan variabel di kiri, konstanta di kanan:
$4x-2x=5+7\Rightarrow 2x=12$.
\langkah{2} $x=6$.
\jadi $x=6$.

\begin{jebakan}
\begin{itemize}[leftmargin=*,topsep=2pt]
  \item \textbf{Salah tanda saat memindah ruas.} Suku yang pindah ruas berganti
  tanda.
\end{itemize}
\end{jebakan}

\begin{variasi}
Persamaan dengan pecahan (kali penyebut dulu); atau yang berujung $0=5$ (tanpa
solusi) / $0=0$ (semua $x$).
\end{variasi}

% ---------------- SOAL 5.3.8 ----------------
\soalno{5.3.8}{pertidaksamaan (balik tanda)}
\begin{soalbox}{}
Selesaikan pertidaksamaan $-2x+1>7$.
\end{soalbox}

\begin{ingat}{Bagi/kali negatif membalik tanda}
Aturan wajib: mengali atau membagi kedua ruas pertidaksamaan dengan bilangan
\emph{negatif} membalik arah tanda ($>$ menjadi $<$, dan sebaliknya).
\end{ingat}

\jawab
\langkah{1} Kurangi $1$: $-2x>6$.
\langkah{2} Bagi $-2$ (\textbf{balik tanda!}): $x<-3$.
\jadi HP $=\{x\mid x<-3\}$.

\begin{jebakan}
\begin{itemize}[leftmargin=*,topsep=2pt]
  \item \textbf{Lupa membalik tanda} sehingga menulis $x>-3$. Ini kesalahan
  nomor satu di pertidaksamaan.
\end{itemize}
\end{jebakan}

\begin{variasi}
$3-5x\le18$; pertidaksamaan ganda $-1<2x+1\le5$ (kerjakan tiap ruas).
\end{variasi}

% ---------------- SOAL 5.3.9 ----------------
\soalno{5.3.9}{membaca makna $m$ dan $c$}
\begin{soalbox}{}
Tarif parkir $y=2000+1000t$ ($t$ jam). Berapa biaya $3$ jam, dan apa makna
$2000$?
\end{soalbox}

\jawab
\langkah{1} Substitusi $t=3$: $y=2000+1000(3)=5000$.
\langkah{2} Tafsir konstanta: $2000$ adalah nilai saat $t=0$, yakni \emph{tarif
awal/masuk} (biaya tetap yang sudah dikenakan meski belum ada jam berjalan);
$1000$ = tarif per jam (gradien).
\jadi biaya $3$ jam $=$ Rp$5.000$; $2000$ = tarif dasar parkir.

\begin{jebakan}
\begin{itemize}[leftmargin=*,topsep=2pt]
  \item \textbf{Mengira $2000$ = tarif per jam.} Yang per jam adalah gradien
  $1000$; $2000$ adalah biaya tetap awal ($t=0$).
\end{itemize}
\end{jebakan}

\begin{variasi}
Tarif taksi $y=7000+4000x$; tagihan listrik dengan abonemen + pemakaian.
\end{variasi}

% ---------------- SOAL 5.3.10 ----------------
\soalno{5.3.10}{model linear terbalik}
\begin{soalbox}{}
Sebuah ojek online menetapkan $y=8000+2500x$. Tentukan jarak $x$ agar ongkos
tepat Rp$23.000$.
\end{soalbox}

\jawab
\langkah{1} Set $y=23000$: $8000+2500x=23000$.
\langkah{2} $2500x=15000\Rightarrow x=6$.
\jadi jaraknya $6$ km.

\begin{jebakan}
\begin{itemize}[leftmargin=*,topsep=2pt]
  \item \textbf{Lupa mengurangi tarif dasar $8000$ dulu.} Bagi $2500$ hanya
  setelah menyisakan bagian yang bergantung jarak.
\end{itemize}
\end{jebakan}

\begin{variasi}
Diberi ongkos, cari jarak maksimum dengan anggaran tertentu (pertidaksamaan
$8000+2500x\le \text{anggaran}$).
\end{variasi}

% ############################################################
\section*{Subbab 5.4 — Program Linear Lanjutan (Pengayaan)}
\addcontentsline{toc}{section}{Subbab 5.4 — Program Linear Lanjutan}
% ############################################################

\begin{ingat}{Langkah baku program linear}
(1) Terjemahkan ke variabel, fungsi objektif $f=ax+by$, dan kendala. (2) Gambar
daerah feasible, tentukan \emph{semua} titik pojok. (3) Hitung $f$ di tiap pojok.
(4) Pilih terbesar (maks) / terkecil (min). \textbf{Optimum selalu di titik
pojok}, tak pernah di dalam daerah.
\end{ingat}

% ---------------- SOAL 5.4.1 ----------------
\soalno{5.4.1}{titik pojok segitiga}
\begin{soalbox}{}
Tentukan titik-titik pojok daerah $x+y\le6$, $x\ge0$, $y\ge0$.
\end{soalbox}

\jawab
\langkah{1} Batas: $x+y=6$ dan kedua sumbu.
\langkah{2} Perpotongan: $(0,0)$; $x+y=6$ dengan sumbu-$x$ $\Rightarrow(6,0)$;
dengan sumbu-$y$ $\Rightarrow(0,6)$.
\jadi titik pojoknya $(0,0),\ (6,0),\ (0,6)$.

\begin{jebakan}
\begin{itemize}[leftmargin=*,topsep=2pt]
  \item \textbf{Salah titik potong sumbu} bila koefisien bukan 1 (mis.\
  $2x+y=6\Rightarrow(3,0),(0,6)$).
\end{itemize}
\end{jebakan}

\begin{variasi}
Tambah kendala membuat pojok bertambah; daerah tak terbatas (arah ``$\ge$'').
\end{variasi}

% ---------------- SOAL 5.4.2 ----------------
\soalno{5.4.2}{maksimum sederhana}
\begin{soalbox}{}
Maksimumkan $f=3x+2y$ pada $x+y\le5$, $x,y\ge0$.
\end{soalbox}

\jawab
\langkah{1} Titik pojok: $(0,0),(5,0),(0,5)$.
\langkah{2} Hitung $f$: $f(0,0)=0$; $f(5,0)=15$; $f(0,5)=10$.
\jadi maksimum $f=15$ di $(5,0)$.

\begin{tips}
Karena koefisien $x$ ($3$) lebih besar dari koefisien $y$ ($2$), masuk akal
optimum condong ke arah memperbanyak $x$ --- yakni pojok $(5,0)$.
\end{tips}

\begin{jebakan}
\begin{itemize}[leftmargin=*,topsep=2pt]
  \item \textbf{Mengambil titik yang ``terjauh''} secara visual tanpa menghitung
  $f$. Selalu evaluasi $f$ di semua pojok.
\end{itemize}
\end{jebakan}

\begin{variasi}
$f=2x+3y$ pada daerah sama memindah optimum ke $(0,5)$.
\end{variasi}

% ---------------- SOAL 5.4.3 ----------------
\soalno{5.4.3}{minimum dengan batas $x\le5$}
\begin{soalbox}{}
Minimumkan $f=x+2y$ pada $x+y\ge4$, $x\le5$, $x,y\ge0$.
\end{soalbox}

\jawab
\langkah{1} Titik pojok relevan pada garis $x+y=4$ dan batas: $(4,0)$; $(0,4)$;
$(5,0)$ (di garis $x=5$, $y=0$, memenuhi $x+y=5\ge4$).
\langkah{2} Hitung $f$: $f(4,0)=4$; $f(0,4)=8$; $f(5,0)=5$.
\jadi minimum $f=4$ di $(4,0)$.

\begin{jebakan}
\begin{itemize}[leftmargin=*,topsep=2pt]
  \item \textbf{Memeriksa titik tak feasible.} Pastikan tiap kandidat pojok
  memenuhi \emph{semua} kendala sebelum menghitung $f$.
\end{itemize}
\end{jebakan}

\begin{variasi}
Tambah kendala atas ($y\le3$) mengubah pojok; ubah objektif ke maksimum (tak ada,
karena daerah membuka ke atas).
\end{variasi}

% ---------------- SOAL 5.4.4 ----------------
\soalno{5.4.4}{perpotongan dua garis}
\begin{soalbox}{}
Tentukan perpotongan garis $x+y=8$ dan $2x+y=12$.
\end{soalbox}

\jawab
\langkah{1} Eliminasi $y$: kurangkan $(2x+y)-(x+y)=12-8\Rightarrow x=4$.
\langkah{2} Substitusi ke $x+y=8$: $4+y=8\Rightarrow y=4$.
\jadi titik potongnya $(4,4)$.

\begin{jebakan}
\begin{itemize}[leftmargin=*,topsep=2pt]
  \item \textbf{Salah tanda saat mengurangkan.} $(2x+y)-(x+y)=x$; kurangi
  seluruh ruas.
\end{itemize}
\end{jebakan}

\begin{variasi}
Perpotongan garis yang butuh perkalian dulu (koefisien tak sama), atau tiga garis
melalui satu titik.
\end{variasi}

% ---------------- SOAL 5.4.5 ----------------
\soalno{5.4.5}{maksimum dengan pojok perpotongan}
\begin{soalbox}{}
Maksimumkan $f=5x+4y$ pada $x+2y\le10$, $3x+y\le15$, $x,y\ge0$.
\end{soalbox}

\jawab
\langkah{1} Titik pojok: $(0,0)$; $(5,0)$ (dari $3x+y=15$ saat $y=0$); $(0,5)$
(dari $x+2y=10$ saat $x=0$); dan perpotongan dua garis miring.
\langkah{2} Perpotongan $x+2y=10$ dan $3x+y=15$: dari kedua, $y=15-3x$; sub:
$x+2(15-3x)=10\Rightarrow x+30-6x=10\Rightarrow -5x=-20\Rightarrow x=4$, lalu
$y=15-12=3$. Titik $(4,3)$.
\langkah{3} Hitung $f=5x+4y$:
\[
f(0,0)=0,\ f(5,0)=25,\ f(0,5)=20,\ f(4,3)=20+12=32.
\]
\jadi maksimum $f=32$ di titik $(4,3)$.

\begin{jebakan}
\begin{itemize}[leftmargin=*,topsep=2pt]
  \item \textbf{Melewatkan pojok perpotongan $(4,3)$.} Justru sering di sanalah
  optimum. Selalu cari titik potong dua kendala miring.
\end{itemize}
\end{jebakan}

\begin{variasi}
Ganti objektif; tambah kendala ketiga sehingga pojok bertambah.
\end{variasi}

% ---------------- SOAL 5.4.6 ----------------
\soalno{5.4.6}{daerah feasible \& pojok}
\begin{soalbox}{}
Gambarkan daerah feasible $2x+y\le8$, $x+3y\le9$, $x,y\ge0$ dan sebutkan titik
pojoknya.
\end{soalbox}

\jawab
\langkah{1} Titik potong sumbu: $2x+y=8\Rightarrow(4,0),(0,8)$;
$x+3y=9\Rightarrow(9,0),(0,3)$.
\langkah{2} Perpotongan dua garis: $y=8-2x$; sub ke $x+3(8-2x)=9\Rightarrow
x+24-6x=9\Rightarrow -5x=-15\Rightarrow x=3$, $y=2$. Titik $(3,2)$.
\langkah{3} Pilih pojok yang \emph{feasible} (memenuhi kedua kendala): $(0,0)$;
$(4,0)$; $(3,2)$; $(0,3)$. (Titik $(0,8)$ dan $(9,0)$ melanggar kendala lain.)

\begin{center}
\begin{tikzpicture}[scale=0.5]
\fill[hijaumuda](0,0)--(4,0)--(3,2)--(0,3)--cycle;
\draw[->](-0.5,0)--(9.5,0) node[right]{$x$};
\draw[->](0,-0.5)--(0,8.8) node[above]{$y$};
\draw[very thick,hijau](4,0)--(0,8);
\draw[very thick,oranye](9,0)--(0,3);
\node at (1.4,1){\small DP};
\fill (3,2) circle(3pt) node[above right]{$(3,2)$};
\end{tikzpicture}
\end{center}
\jadi titik pojoknya $(0,0),(4,0),(3,2),(0,3)$.

\begin{jebakan}
\begin{itemize}[leftmargin=*,topsep=2pt]
  \item \textbf{Memasukkan titik potong sumbu yang tak feasible} ($(0,8)$
  melanggar $x+3y\le9$). Uji tiap kandidat.
\end{itemize}
\end{jebakan}

\begin{variasi}
Tambahkan fungsi objektif dan cari optimum di pojok-pojok ini.
\end{variasi}

% ---------------- SOAL 5.4.7 ----------------
\soalno{5.4.7}{model \& maksimasi laba penjahit}
\begin{soalbox}{}
Seorang penjahit membuat baju ($x$) dan celana ($y$); baju butuh $2$ m kain,
celana $3$ m, tersedia $24$ m. Laba baju Rp$40$rb, celana Rp$50$rb, dan maksimal
$10$ baju. Modelkan dan maksimalkan laba.
\end{soalbox}

\jawab
\langkah{1} Kendala: $2x+3y\le24$ (kain), $x\le10$ (batas baju), $x,y\ge0$.
Objektif (ribuan): $f=40x+50y$.
\langkah{2} Titik pojok daerah:
\begin{itemize}[leftmargin=*,topsep=1pt]
  \item $(0,0)$;\quad $(10,0)$ ($x=10,y=0$; $2\cdot10=20\le24$ \checkmark);
  \item $(10,\tfrac43)$ (dari $x=10$ dan $2x+3y=24\Rightarrow 3y=4$);
  \item $(0,8)$ (dari $x=0$, $3y=24$).
\end{itemize}
\langkah{3} Hitung $f=40x+50y$:
\[
f(0,0)=0,\ f(10,0)=400,\ f(0,8)=400,\ f\!\left(10,\tfrac43\right)=400+\tfrac{200}{3}\approx466{,}7.
\]
\jadi laba maksimum $\approx$ Rp$466.667$ di $(10,\tfrac43)$. Karena celana harus
bulat, dalam praktik ambil $x=10,\ y=1$ (laba Rp$450$rb) sebagai solusi bilangan
bulat terdekat yang feasible.

\begin{jebakan}
\begin{itemize}[leftmargin=*,topsep=2pt]
  \item \textbf{Lupa kendala $x\le10$.} Tanpa itu, pojok berpindah dan jawaban
  berubah. Baca semua batasan.
  \item \textbf{Menganggap solusi LP selalu bulat.} Optimum matematis bisa
  pecahan; keputusan produksi lalu dibulatkan.
\end{itemize}
\end{jebakan}

\begin{variasi}
Tambah kendala jam kerja; ubah laba per unit sehingga optimum berpindah pojok.
\end{variasi}

% ---------------- SOAL 5.4.8 ----------------
\soalno{5.4.8}{mengapa optimum di pojok}
\begin{soalbox}{}
Mengapa optimum fungsi linear tidak mungkin berada di bagian dalam daerah
feasible? Jelaskan singkat.
\end{soalbox}

\jawab
\langkah{1} Perhatikan ``garis selidik'' $ax+by=k$: himpunan titik dengan nilai
objektif sama. Mengubah $k$ menggeser garis ini \emph{sejajar}.
\langkah{2} Dari titik mana pun di \emph{dalam} daerah, kita masih bisa bergerak
searah yang memperbesar (atau memperkecil) $f$ tanpa keluar daerah --- jadi titik
dalam \emph{tidak pernah} ekstrem.
\langkah{3} Gerakan itu baru berhenti saat garis selidik menyentuh daerah
terakhir kali, yaitu tepat di sebuah \emph{titik pojok} (atau sepanjang sisi).
\jadi optimum selalu tercapai di titik pojok (teorema titik pojok).

\begin{jebakan}
\begin{itemize}[leftmargin=*,topsep=2pt]
  \item \textbf{Mengira ada ``puncak'' di tengah.} Fungsi linear tak punya bukit
  atau lembah di interior --- ia datar-miring, jadi ekstrem hanya di tepi.
\end{itemize}
\end{jebakan}

\begin{variasi}
Kasus objektif \emph{sejajar} salah satu sisi $\Rightarrow$ optimum di seluruh
sisi (tak tunggal), tetapi tetap mencakup titik pojok.
\end{variasi}

% ---------------- SOAL 5.4.9 ----------------
\soalno{5.4.9}{maks \& min pada segitiga}
\begin{soalbox}{}
Pada $f=2x+3y$ di daerah segitiga $(0,0),(4,0),(0,3)$, tentukan maksimum dan
minimumnya.
\end{soalbox}

\jawab
\langkah{1} Evaluasi $f$ di ketiga pojok:
\[
f(0,0)=0,\quad f(4,0)=8,\quad f(0,3)=9.
\]
\langkah{2} Bandingkan.
\jadi \textbf{maksimum} $f=9$ di $(0,3)$; \textbf{minimum} $f=0$ di $(0,0)$.

\begin{jebakan}
\begin{itemize}[leftmargin=*,topsep=2pt]
  \item \textbf{Melupakan bahwa min \& maks keduanya di pojok.} Cukup bandingkan
  nilai di semua pojok; yang terbesar maks, terkecil min.
\end{itemize}
\end{jebakan}

\begin{variasi}
Objektif $f=x-y$ bisa memberi min negatif; daerah dengan pojok perpotongan.
\end{variasi}

% ---------------- SOAL 5.4.10 ----------------
\soalno{5.4.10}{model minimasi biaya pupuk}
\begin{soalbox}{}
Suatu kebun butuh pupuk minimal $12$ kg N dan $8$ kg P. Pupuk A (per karung)
memberi $2$ N, $1$ P seharga Rp$50$rb; pupuk B memberi $1$ N, $2$ P seharga
Rp$40$rb. Minimalkan biaya (rumuskan kendala dan objektifnya, lalu selesaikan).
\end{soalbox}

\jawab
\langkah{1} Misalkan $x=$ karung A, $y=$ karung B. Kendala nutrisi (minimal
$\Rightarrow$ ``$\ge$''):
\[
2x+y\ge12\ (\text{N}),\qquad x+2y\ge8\ (\text{P}),\qquad x,y\ge0.
\]
Objektif (minimasi, ribuan): $f=50x+40y$.
\langkah{2} Titik pojok feasible: $(0,12)$ (dari $2x+y=12$, $x=0$); $(8,0)$
(dari $x+2y=8$, $y=0$)? cek N: $2\cdot8=16\ge12$ \checkmark; dan perpotongan
$2x+y=12$ dengan $x+2y=8$.
\langkah{3} Perpotongan: dari $y=12-2x$; sub $x+2(12-2x)=8\Rightarrow
x+24-4x=8\Rightarrow -3x=-16\Rightarrow x=\tfrac{16}{3}$, $y=12-\tfrac{32}{3}
=\tfrac{4}{3}$.
\langkah{4} Hitung $f$:
\[
f(0,12)=480,\quad f(8,0)=400,\quad f\!\left(\tfrac{16}{3},\tfrac43\right)
=50\cdot\tfrac{16}{3}+40\cdot\tfrac43=\tfrac{800+160}{3}=\tfrac{960}{3}=320.
\]
\jadi biaya minimum Rp$320.000$ di $\left(\tfrac{16}{3},\tfrac43\right)$
(sekitar $5{,}3$ karung A dan $1{,}3$ karung B; dibulatkan sesuai kebijakan
pembelian).

\begin{jebakan}
\begin{itemize}[leftmargin=*,topsep=2pt]
  \item \textbf{Memakai ``$\le$'' untuk kebutuhan minimal.} ``Minimal $12$ kg''
  berarti $\ge12$, bukan $\le$.
  \item \textbf{Berhenti di titik sumbu.} Sering minimum ada di titik potong dua
  kendala, bukan di sumbu.
\end{itemize}
\end{jebakan}

\begin{variasi}
Tambah nutrisi ketiga (K) sebagai kendala; ubah harga sehingga optimum berpindah.
\end{variasi}

% ============================================================
\section*{Jawaban Studi Kasus Pengayaan (Program Linear)}
\addcontentsline{toc}{section}{Jawaban Studi Kasus Pengayaan}

\begin{studikasus}{Pabrik Mebel}
Meja ($x$): $4$ jam tukang, $2$ kg kayu; kursi ($y$): $2$ jam, $3$ kg. Tersedia
$40$ jam dan $42$ kg kayu. Laba meja Rp$80$rb, kursi Rp$60$rb.
\end{studikasus}
\textbf{Jawaban.} Kendala $2x+y\le20$ (dari $4x+2y\le40$), $2x+3y\le42$,
$x,y\ge0$. Objektif $f=80x+60y$ (ribuan). Pojok: $(0,0),(10,0),(0,14)$, dan
perpotongan $2x+y=20$ \& $2x+3y=42$ $\Rightarrow$ kurangkan: $2y=22\Rightarrow
y=11$, $x=4{,}5$. Uji: $f(10,0)=800$, $f(0,14)=840$, $f(4{,}5;11)=360+660=1020$.
\jadi laba maksimum \textbf{Rp$1.020.000$} di $(4{,}5;\ 11)$ (dibulatkan sesuai
kebijakan produksi).

\begin{studikasus}{Pengiriman Barang}
Truk kecil ($x$): $2$ ton, Rp$300$rb/rit; truk besar ($y$): $5$ ton,
Rp$600$rb/rit. Kirim minimal $30$ ton; total rit maksimal $10$. Minimumkan ongkos.
\end{studikasus}
\textbf{Jawaban.} Kendala $2x+5y\ge30$, $x+y\le10$, $x,y\ge0$. Objektif
$f=300x+600y$ (ribuan). Pojok feasible meliputi $(0,6)$ (dari $5y=30$),
$(0,10)$, dan perpotongan $2x+5y=30$ \& $x+y=10$ $\Rightarrow x=\tfrac{20}{3},
y=\tfrac{10}{3}$. Uji: $f(0,6)=3600$; $f(0,10)=6000$;
$f(\tfrac{20}{3},\tfrac{10}{3})=2000+2000=4000$.
\jadi ongkos minimum \textbf{Rp$3.600.000$} di $(0,6)$ --- gunakan $6$ truk besar.

% ============================================================
\begin{tips}
\textbf{Ringkasan strategi Bab 5.}
\begin{enumerate}[leftmargin=*,topsep=2pt]
  \item \emph{SPLTV:} eliminasi bertingkat ($3\to2\to1$ variabel); jawaban boleh
  pecahan; selalu verifikasi dengan substitusi.
  \item \emph{Banyak solusi:} berpotongan (tunggal), berimpit (tak hingga),
  sejajar (nihil).
  \item \emph{SPtLDV:} gambar garis, \textbf{uji titik} untuk arah arsiran,
  ambil irisan. Jangan lupa $x,y\ge0$ pada konteks nyata.
  \item \emph{Program linear:} optimum di \textbf{titik pojok}. Cari semua pojok
  (termasuk perpotongan kendala miring), hitung $f$, pilih terbaik.
\end{enumerate}
\end{tips}

\end{document}
